\chapter{Discussion}
\label{09_diskussion}

This thesis developed and evaluated a proof-of-concept workflow for improving existing building footprint polygons in selected study areas rather than generating a completely new building dataset from scratch. This distinction is important for geospatial data maintenance. In many practical settings, especially where automatically generated building datasets already exist, the central challenge is not only to detect buildings again, but to assess whether available polygons are fit for use and to selectively correct them where the imagery indicates geometric errors. This follows the broader idea of geodatabase updating, where modifying and improving existing spatial data can be more efficient than repeatedly creating new datasets \cite{Qin_changedetection}.

The work therefore connects two research directions that are often treated separately. On the one hand, building footprint extraction from high-resolution imagery has been studied extensively using classical deep learning architectures such as patch-based CNNs, fully convolutional networks, SegNet, and U-Net variants \cite{Luo_overviewbuioldingfootprint, rastogi2022automatic}. On the other hand, recent research on foundation models and multimodal large language models suggests that general-purpose models can support visual reasoning, spatial interpretation, and task-specific prompting without being trained from scratch for every individual application \cite{bommasani2022opportunitiesrisksfoundationmodels, wu2023multimodallargelanguagemodels, Yin_2024MLLM}. In this thesis, these ideas are combined in a two-stage workflow: a multimodal model first evaluates the plausibility of an existing footprint in relation to the image context, and a segmentation foundation model then refines the geometry where correction appears meaningful.

The specific contribution of this approach lies in the separation between semantic quality assessment and geometric correction. The MLLM is used for tasks that require contextual interpretation, such as deciding whether a visible roof is present, whether the original polygon covers the complete structure, and which type of error is likely present. SAM is then used for the image-driven segmentation task, where its prompt-based design allows the model to transform the MLLM output into a refined mask \cite{kirillov2023segment}. The final polygon is not simply the raw segmentation result, but the outcome of a controlled post-processing workflow that regularises and simplifies the segmentation into a usable vector geometry.

This setup is particularly relevant for data-scarce regions in the Global South. Global building footprint products such as Google Open Buildings provide important spatial information where authoritative datasets are often incomplete or unavailable \cite{sirko2021continentalscalebuildingdetectionhigh, CHAMBERLAIN2024102104}. At the same time, such automatically derived products contain characteristic errors related to shifts, incomplete footprints, merged structures, vegetation occlusion, and settlement-specific roof patterns. The discussion therefore focuses on whether the proposed combination of MLLM-based quality assessment, SAM-based refinement, and deterministic post-processing can meaningfully improve these existing polygons, where it fails, and what the results imply for the transferability of foundation-model-based geospatial quality assurance across different study areas.




\newpage


